\subsection{桶与桶型空间}

命题 1. --- 设 T 是 E 的一个子集。下列条件等价：

\begin{enumerate}
\def\labelenumi{(\roman{enumi})}
\item
  T 是凸的、均衡的、闭的且吸收的。
\item
  T 是 E' 中一个凸的、均衡的且弱有界的集 M 的极。
\item
  存在一个下半连续半范数 \(p\)（在 \(\mathbf{E}\) 上），使得 \(\mathbf{T}\) 是所有 \(x \in \mathbf{E}\) 中使 \(p(x) \leqslant 1\) 者所成之集。
\item
  \(\Rightarrow\) (ii)：在 (i) 的假设下，设 \(\mathbf{M} = \mathbf{T}^{\circ}\)；则 \(\mathbf{M}\) 在 \(\mathbf{E}'\) 中是凸的且均衡的。对每个 \(x \in \mathbf{E}\)，存在实数 \(r > 0\) 使得 \(rx \in \mathbf{T}\)，因而 \(|u(x)| \leqslant \frac{1}{r}\) 对所有 \(u \in \mathbf{M}\) 成立；换言之，\(\mathbf{M}\) 是弱有界的。由 II, 第 45 页 的推论 3，我们有 \(\mathrm{T} = \mathbf{M}^{\circ}\)，故 \(\mathrm{T}\) 满足 (ii)。
\item
  \(\Rightarrow\) (iii)：在 (ii) 的假设下，设 \(p(x) = \sup_{u\in \mathbf{M}}|u(x)|\)（对所有 \(x\in \mathrm{E}\)）。立即可见 \(T = M^{\circ}\) 由所有 \(x \in E\) 中使 \(p(x) \leqslant 1\) 者组成。\(E'\) 上的半范数 p 是下半连续的，因为它是一族连续函数的上包络（GT, IV, § 6, No.~2, 推论）。
\item
  \(\Rightarrow\) (i)：这是显然的。
\end{enumerate}

推论. --- 下列条件等价：

\begin{enumerate}
\def\labelenumi{(\roman{enumi})}
\item
  \(\mathrm{E}'\) 的每个弱有界子集都是等度连续的；
\item
  E 中每个凸的、均衡的、闭的且吸收的集都是 0 的一个邻域；
\item
  E 上每个下半连续半范数都是连续的。
\end{enumerate}

定义 1. --- 称满足命题 1 的诸等价条件的集 T 为 E 中的一个桶。

定义 2. --- 称空间 E 为桶型的，如果它满足命题 1 的推论的诸等价条件。

我们知道（III, 第 22 页, 命题 9），E 的对偶 E' 的每个等度连续子集都是强有界且弱有界的。因此我们可以把桶型空间的定义重述如下：

评注. --- 在桶型空间的对偶中，等度连续集的类、强有界集的类、弱有界集的类以及弱拓扑下相对紧集的类全都是相同的。若 E 是 Hausdorff 且桶型的，且 \(E_{b}'\) 是它的强对偶，则其中一个空间的 0 的诸邻域的极构成另一个空间的典范有界结构的一个基，而其中一个空间的有界子集的极构成另一个空间的 0 的邻域滤子的一个基。

命题 2. --- 每个是 Baire 空间（GT, IX, § 5, No.~3）的局部凸空间 E 都是桶型的。

设 \(\mathrm{T}\) 是 \(\mathrm{E}\) 中的一个桶；由于 \(\mathrm{T}\) 是吸收的，\(\mathrm{E}\) 是诸闭集 \(n\mathrm{T}\)（\(n\) 为整数 \(> 0\)）的并；由于 \(\mathrm{E}\) 是 Baire 空间，这些集中至少有一个含有内点，因而 T 本身有一个内点 x。若 \(x \neq 0\)，则由于 \(-x \in T\)，且 0 是以 x 和 -x 为端点的开线段上的一点，0 是凸集 T 的一个内点（II, 第 14 页, 命题 16）。因此 T 是 0 的一个邻域。

推论. --- 每个 Fréchet 空间（特别地，每个 Banach 空间）都是桶型的。这由 Baire 定理（GT, IX, § 5, No.~3, 定理 1）得出。

命题 3. --- 设 \((\mathrm{F}_{i})_{i\in\mathrm{I}}\) 是桶型空间的一个族，且对每个 \(i\in I\)，设 \(f_{i}\) 是从 \(F_{i}\) 到一个向量空间 E 中的一个线性映射。带使诸 \(f_{i}\) 连续的最细局部凸拓扑（II, 第 27 页, 命题 5）的空间 E 是一个桶型空间。

设 T 是 E 中的一个桶。由于 \(f_{i}\) 连续，\(f_{i}^{-1}(T)\) 是 \(F_{i}\) 中一个凸的、均衡的、闭的且吸收的集；换言之，是 \(F_{i}\) 中的一个桶；由于 \(F_{i}\) 是桶型的，\(f_{i}^{-1}(T)\) 是 \(F_{i}\) 中 0 的一个邻域（对所有 \(i \in I\)）。这蕴含 T 是 E 中 0 的一个邻域（II, 第 27 页, 命题 5）。

推论 1. --- 桶型空间的每个商空间都是桶型的。

推论 2. --- 设 \((\mathrm{E}_{i})_{i\in\mathrm{I}}\) 是局部凸空间的一个族，并设 E 是这一族的拓扑直和。为使 E 是桶型的，必要且充分的是每个 \(E_{i}\) 都是桶型的。

由命题 3，这一条件显然是充分的；它也是必要的，由推论 1，因为每个 \(E_{i}\) 同构于 E 的一个商空间（II, 第 31 页, 命题 8）。

推论 3. --- 桶型空间的每个归纳极限都是桶型空间。

我们将在后面证明（IV, 第 14 页, 推论）桶型空间的每个乘积都是桶型的。

